\documentclass[12pt,a4paper]{article}

%% Kodowanie i język polski
\usepackage{fontspec}
\usepackage{polyglossia}
\setmainlanguage{polish}

%% Czcionki – wymagania ATINS (TNR → Liberation Serif – metryczny zamiennik)
\setmainfont{Liberation Serif}

%% Geometria strony
\usepackage{geometry}
\geometry{
  a4paper,
  left=2.5cm,
  right=2.5cm,
  top=2.5cm,
  bottom=2.5cm,
}

%% Interlinia
\usepackage{setspace}
\onehalfspacing

%% Przypisy dolne – 10 pkt (wymaganie ATINS)
\usepackage[bottom]{footmisc}
\renewcommand{\footnotesize}{\fontsize{10pt}{12pt}\selectfont}

%% Pakiety naukowe
\usepackage{graphicx}
\usepackage{booktabs}
\usepackage{amsmath}
\usepackage{csquotes}
\usepackage[hidelinks]{hyperref}
\usepackage{xcolor}
\usepackage{listings}
\usepackage{enumitem}
\usepackage{caption}
\usepackage{tikz}
\usetikzlibrary{shapes.geometric, arrows.meta, positioning, fit, backgrounds, shadows}

%% Listingi kodu
\definecolor{codegreen}{rgb}{0,0.6,0}
\definecolor{codegray}{rgb}{0.5,0.5,0.5}
\definecolor{codepurple}{rgb}{0.58,0,0.82}
\definecolor{backcolour}{rgb}{0.96,0.96,0.94}

\lstdefinestyle{mystyle}{
    backgroundcolor=\color{backcolour},
    commentstyle=\color{codegreen},
    keywordstyle=\color{magenta}\bfseries,
    numberstyle=\tiny\color{codegray},
    stringstyle=\color{codepurple},
    basicstyle=\linespread{1.0}\ttfamily\footnotesize\selectfont,
    breakatwhitespace=false,
    breaklines=true,
    captionpos=b,
    keepspaces=true,
    numbers=left,
    numbersep=5pt,
    showspaces=false,
    showstringspaces=false,
    showtabs=false,
    tabsize=2
}
\lstset{style=mystyle}

%% Bibliografia
\usepackage[
  backend=biber,
  style=numeric,
  sorting=none,
  maxnames=3,
  minnames=1,
]{biblatex}
\addbibresource{bibliography.bib}

%%----------------------------------------------
\begin{document}
\sloppy

%% Tytuł i autor
\title{%
  \Large\bfseries
  Bezpieczna kompartmentaryzacja stacji roboczej\\
  w środowisku akademickim: wykorzystanie Qubes OS\\
  w działaniach OSINT, programowaniu oraz bezpiecznej detonacji i weryfikacji prac studenckich%
}
\author{Przemysław Żarnecki\\
  \small Akademia Techniczno-Informatyczna Nauk Stosowanych we Wrocławiu\\
  \small \texttt{pzarnecki@protonmail.com}
}
\date{}
\maketitle

%% Streszczenie
\begin{abstract}
\noindent
Współczesne środowisko akademickie stawia przed studentami i nauczycielami akademickimi wyzwania wymagające operowania na niezweryfikowanych danych, repozytoriach kodu oraz zasobach sieciowych. Działania takie jak pozyskiwanie informacji jawnoźródłowych (OSINT), analiza ofert rynkowych, testowanie zewnętrznych bibliotek programistycznych podatnych na ataki łańcucha dostaw oraz laboratoryjna ocena nadsyłanego przez studentów kodu źródłowego niosą wysokie ryzyko naruszenia prywatności, profilowania cyfrowego lub kompromitacji stacji roboczej. Niniejszy artykuł przedstawia kompleksowy model bezpiecznej stacji akademickiej opartej na kompartmentaryzacji sprzętowej z wykorzystaniem hipernadzorcy typu 1 (Xen) i systemu Qubes OS 4.2. W pracy zdefiniowano formalny model stref zaufania (Academic Compartmentalization Framework -- ACF), zaproponowano bezśladowy potok wywiadu jawnoźródłowego oparty na jednorazowych maszynach Whonix oraz opracowano mechanizm Zero-Trust Student Code Evaluation Pipeline (ZT-SCEP), wykorzystujący Qubes RPC (\texttt{qrexec}) do zautomatyzowanej detonacji i oceny prac laboratoryjnych bez dostępu do sieci. Przeprowadzona ewaluacja empiryczna wykazała, że narzut pamięciowy i czas uruchamiania środowisk efemerycznych (średnio 2,4~s) są w pełni akceptowalne w codziennej pracy dydaktycznej, zapewniając bezwzględną odporność stacji bazowej na próby ucieczki kodu, pętle widelcowe (\textit{fork bombs}) i eksfiltrację tożsamości kryptograficznej.

\medskip
\noindent
\textbf{Słowa kluczowe:} Qubes OS, kompartmentaryzacja, OSINT, cyberbezpieczeństwo akademickie, detonacja kodu, qrexec, ocena prac studenckich, ochrona prywatności, suwerenność cyfrowa.
\end{abstract}

%%----------------------------------------------
\section{Wprowadzenie}
\label{sec:wprowadzenie}

Codzienna praca badawczo-dydaktyczna w uczelniach technicznych wymaga przetwarzania danych z niepewnych źródeł zewnętrznych. Wykładowcy, naukowcy i studenci stają w obliczu trzech kluczowych wektorów operacyjnych:
\begin{enumerate}[label=\alph*)]
    \item \textbf{Wywiad jawnoźródłowy (OSINT) i monitoring rynku:} poszukiwanie literatury, analiza repozytoriów, weryfikacja ofert praktyk i pracy w portalach rekrutacyjnych oraz monitoring cyberzagrożeń. Działania te narażają na agresywne profilowanie behawioralne, deanonymizację tożsamości akademickiej oraz pobieranie spreparowanych plików (np. dokumentów PDF z exploitami);
    \item \textbf{Wytwarzanie oprogramowania a łańcuch dostaw:} laboratoryjne projekty programistyczne opierają się na ekosystemach \texttt{npm}, \texttt{pip} czy \texttt{cargo}. Narastająca liczba ataków typu \textit{supply chain} oraz \textit{typosquatting} sprawia, że instalacja pakietu może skutkować kradzieżą kluczy SSH lub tokenów API z dysku dewelopera~\cite{ohm2020backstabber,ladisa2023sok};
    \item \textbf{Ocena nadsyłanych prac laboratoryjnych:} nauczyciel akademicki otrzymuje w semestrze setki archiwów z kodem źródłowym i binariami. Uruchomienie takiego kodu na stacji prowadzącego niesie ryzyko: od przypadkowych pętli forkowania (\textit{fork bomb}), po złośliwe skrypty eksfiltrujące dane przez sieć uczelnianą~\cite{edwards2017grading}.
\end{enumerate}

Tradycyjne podejście oparte na systemach monolitycznych z antywirusami jest bezsilne wobec ataków ukierunkowanych i nieznanych podatności (\textit{zero-day})~\cite{nask_cyberuczelnia2025,enisa_threat_landscape2024}. Z kolei konteneryzacja (np. Docker) nie zapewnia twardej granicy bezpieczeństwa ze względu na współdzielenie jądra (\textit{shared-kernel}) i częste ucieczki do hosta (\textit{container breakout})~\cite{sultan2019containers}. 

Odpowiedzią na te zagrożenia jest paradygmat \textbf{bezpieczeństwa przez kompartmentaryzację} (\textit{security by compartmentalization}), zaimplementowany w systemie \textbf{Qubes OS}~\cite{rutkowska2010qubes}. Wykorzystując hipernadzorca typu 1 (Xen), Qubes OS izoluje poszczególne sfery aktywności w niezależnych, sprzętowo odizolowanych maszynach wirtualnych (AppVMs i Disposable VMs).

Celem niniejszego artykułu jest opracowanie, wdrożenie i empiryczna ewaluacja zintegrowanego modelu stacji akademickiej opartej na systemie Qubes OS 4.2. W pracy przedstawiono nowatorskie mechanizmy:
\begin{itemize}
    \item koncepcję zanonimizowanego i bezśladowego potoku badawczego OSINT opartego na jednorazowych maszynach Whonix, eliminującego ryzyko profilowania tożsamości akademickiej;
    \item architekturę separacji tożsamości deweloperskiej z wykorzystaniem mechanizmów Split-SSH i Split-GPG, zabezpieczającą proces badawczo-rozwojowy przed kompromitacją łańcucha dostaw;
    \item autorski potok \textit{Zero-Trust Student Code Evaluation Pipeline} (ZT-SCEP), realizujący bezpieczną, automatyczną detonację i weryfikację prac studenckich przy użyciu magistrali Qubes RPC (\texttt{qrexec}) w efemerycznych domenach bez dostępu do sieci.
\end{itemize}

\section{Przegląd literatury i stan badań}
\label{sec:przeglad_literatury}

W literaturze przedmiotu paradygmat cyberbezpieczeństwa ulega fundamentalnej redefinicji. Konwencjonalne podejście \textit{detect and block} (oparte na sygnaturach antywirusowych i heurystyce) ustępuje miejsca doktrynie \textbf{Assume Breach} (założenie o nieuchronnym włamaniu), której wiodącą realizacją technologiczną jest sprzętowa kompartmentaryzacja stacji roboczej~\cite{rutkowska2010qubes,kapera2026methodology}.

\subsection{Kompartmentaryzacja a tradycyjne modele wirtualizacji}
Podstawy architektury Qubes OS sformułowali Rutkowska i Wojtczuk~\cite{rutkowska2010qubes}, wskazując, że monolityczne jądra współczesnych systemów operacyjnych (liczące miliony linii kodu) oraz podatne sterowniki sprzętowe uniemożliwiają zapewnienie integralności pojedynczej instancji systemu. W odróżnieniu od komercyjnych hipernadzorców typu 2 (takich jak VirtualBox czy VMware Workstation), instalowanych na wierzchu podatnego systemu operacyjnego gospodarza, Qubes OS implementuje hypervisor typu 1 (bare-metal) \textbf{Xen}~\cite{barham2003xen}.

W architekturze tej uprzywilejowana domena administracyjna (\texttt{dom0}) jest całkowicie pozbawiona łączności sieciowej. Sterowniki kart sieciowych i magistrali USB są wyizolowane w dedykowanych maszynach domenowych (\texttt{sys-net}, \texttt{sys-usb}), komunikujących się poprzez sprzętowy mechanizm IOMMU (Intel VT-d / AMD-Vi). Jak wykazali w 2026 roku Kapera i Niemiec~\cite{kapera2026methodology} w kompleksowych badaniach empirycznych opartych na macierzy MITRE ATT\&CK, architektura ta redukuje przestrzeń podatności stacji roboczej o ponad 80\% w porównaniu z wiodącym systemem Windows 11. Kluczowym wyróżnikiem jest uniemożliwienie trwałej infekcji (\textit{persistence}) oraz przemieszczania lateralnego (\textit{lateral movement}) poprzez zastosowanie maszyn efemerycznych (\textit{Disposable VMs} -- \texttt{DispVM}), których stan dyskowy jest bezpowrotnie niszczony w momencie zamknięcia procesu~\cite{kapera2026methodology}.

\subsection{Wyzwania operacyjne OSINT i higiena tożsamości cyfrowej}
Prowadzenie badań typu OSINT przez studentów i naukowców wiąże się z ryzykiem zaawansowanego profilowania cyfrowego. Zgodnie z badaniami Bazzella~\cite{bazzell2023osint}, nowoczesne platformy internetowe stosują pasywne i aktywne mechanizmy identyfikacji użytkownika, takie jak fingerprinting silnika renderującego (Canvas, WebGL), profilowanie stosu TCP/IP oraz analizę ciasteczek śledzących. 

Połączenie systemu Qubes OS z dedykowaną dystrybucją Whonix~\cite{whonix_design2024} realizuje rygorystyczny dwuetapowy model izolacji. Maszyna bramy sieciowej (\texttt{sys-whonix}) wymusza tunelowanie całego ruchu przez rozproszoną sieć Tor~\cite{dingledine2004tor}, podczas gdy stacja robocza (\texttt{anon-whonix}) nie posiada wiedzy o fizycznym adresie IP, MAC ani parametrach sprzętowych karty sieciowej. W rezultacie nawet potencjalny exploit przeglądarki klasy RCE (Remote Code Execution) nie pozwala atakującemu na deanonymizację badacza, gdyż nie ma on programistycznego dostępu do stosu sieciowego gospodarza.

\subsection{Bezpieczeństwo łańcucha dostaw w procesie programistycznym}
Dynamiczny rozwój otwartoźródłowych repozytoriów kodu doprowadził do gwałtownego wzrostu ataków na łańcuch dostaw oprogramowania (\textit{software supply chain attacks}). Ohm i wsp.~\cite{ohm2020backstabber} oraz Ladisa i wsp.~\cite{ladisa2023sok} sklasyfikowali wektory infekcji za pośrednictwem menedżerów pakietów (\texttt{npm}, \texttt{pip}, \texttt{crates.io}). Powszechną techniką jest \textit{typosquatting} oraz ukierunkowane wstrzykiwanie złośliwego kodu do skryptów instalacyjnych (np. \texttt{setup.py} lub hooków \texttt{preinstall}). 

Głównym celem takich ataków jest bezszelestna eksfiltracja poświadczeń zgromadzonych na stacji roboczej programisty: kluczy prywatnych SSH (\texttt{\textasciitilde/.ssh}), tokenów API do chmur publicznych oraz kluczy prywatnych GPG. W tradycyjnym systemie jednodomenowym dowolny uruchomiony proces posiada uprawnienia do odczytu całego katalogu domowego zalogowanego użytkownika. Zastosowanie kompartmentaryzacji w standardzie Split-SSH i Split-GPG w Qubes OS pozwala fizycznie odseparować magazyn kluczy kryptograficznych od zanieczyszczonego środowiska kompilacji.

\subsection{Ryzyka związane z automatyczną oceną prac studenckich}
Kwestia bezpieczeństwa weryfikacji prac laboratoryjnych z informatyki była badana m.in. przez Edwardsa i Shams~\cite{edwards2017grading}. Autorzy wykazali, że nadsyłane przez studentów zadania programistyczne często wykazują cechy kodu wrogiego -- zarówno w wyniku celowego działania (np. próby ucieczki z piaskownicy, odczytania szablonów oceniania), jak i rażących błędów implementacyjnych (wycieki pamięci, rekurencyjne pętle \texttt{fork()}, próby otwarcia tysięcy gniazd sieciowych).

Stosowane w uczelniach lekkie rozwiązania wirtualizacji kontenerowej (np. Docker) są powszechnie krytykowane w literaturze za brak twardej bariery izolacyjnej~\cite{sultan2019containers}. Kontenery współdzielą jądro z systemem gospodarza, przez co podatności jądra Linuksa (np. wyścigi uprawnień \textit{Dirty COW}, błędy w podsystemie \textit{io\_uring}) umożliwiają natychmiastowe przejęcie uprawnień administratora (\texttt{root}) na maszynie sprawdzającej.

Dotychczasowe badania nad infrastrukturą uczelnianą koncentrowały się na serwerowych klastrach wirtualizacji (np. zarządzanie środowiskami wirtualnymi i modelami językowymi w Proxmox~VE~\cite{zarnecki2026llm_proxmox}). W literaturze brakuje jednak opracowania, które adaptowałoby architekturę Qubes OS bezpośrednio do wieloaspektowego profilu akademickiej stacji roboczej, integrując w jednym spójnym środowisku bezpieczny OSINT, kompartmentaryzowany development oraz automatyczną detonację kodu studenckiego. Niniejsza praca wypełnia tę istotną lukę badawczą.

\section{Architektura systemu i formalny model zaufania (ACF)}
\label{sec:architektura_i_model}

W celu zapewnienia bezkompromisowego bezpieczeństwa stacji roboczej pracownika naukowo-dydaktycznego oraz studenta, w ramach niniejszych badań opracowano formalny model wielodomenowej separacji -- \textbf{Academic Compartmentalization Framework (ACF)}.

\subsection{Koncepcja stref zaufania i izolacji domenowej}
Model ACF kategoryzuje operacje akademickie na pięć ortogonalnych warstw zaufania, mapowanych bezpośrednio na domeny wirtualne (\textit{qubes}) hipernadzorcy Xen, zgodnie z zasadą najmniejszych uprawnień (\textit{Principle of Least Privilege}):
\begin{enumerate}
    \item \textbf{Warstwa sprzętowa i sterowników (IOMMU Passthrough):}
    \begin{itemize}
        \item \texttt{sys-net} -- domena izolująca fizyczne kontrolery Ethernet i Wi-Fi. Ewentualna kompromitacja sterownika bezprzewodowego (częsty wektor ataków radiowych) nie ma wpływu na system macierzysty;
        \item \texttt{sys-usb} -- domena izolująca kontrolery USB. Skutecznie neutralizuje wektory sprzętowe, takie jak złośliwe kontrolery firmware (\textit{BadUSB}) nadsyłane lub przynoszone przez studentów na nośnikach pamięci flash;
        \item \texttt{sys-firewall} -- pośrednicząca zapora sieciowa filtrująca ruch pomiędzy domenami roboczymi a bramami wyjściowymi.
    \end{itemize}
    \item \textbf{Warstwa tożsamości i kryptografii (\textit{Air-Gapped Vaults}):}
    \begin{itemize}
        \item \texttt{vault-sec} -- całkowicie odcięta od sieci wirtualna maszyna przechowująca prywatne klucze PGP/GPG, certyfikaty tożsamości X.509, klucze ed25519 dla protokołu SSH oraz bazę haseł. Magazyn ten nigdy nie wchodzi w bezpośrednią interakcję ze światem zewnętrznym.
    \end{itemize}
    \item \textbf{Warstwa zaufanej pracy akademickiej (\textit{Trusted Academic Realm}):}
    \begin{itemize}
        \item \texttt{work-academic} -- domena przeznaczona do edycji dokumentów, przygotowywania materiałów dydaktycznych, obsługi poczty instytucjonalnej oraz wprowadzania ocen do uczelnianego systemu dziekanatowego;
        \item \texttt{dev-trusted} -- środowisko deweloperskie zaufanych projektów badawczych i publikacyjnych, posiadające dostęp do sieci zewnętrznej w celach synchronizacji repozytoriów git.
    \end{itemize}
    \item \textbf{Warstwa wywiadu jawnoźródłowego (\textit{OSINT Realm}):}
    \begin{itemize}
        \item \texttt{sys-whonix} i efemeryczne instancje \texttt{disp-anon-whonix} -- izolowane środowisko analizy źródeł otwartych, uniemożliwiające profilowanie i śledzenie adresu IP uczelni.
    \end{itemize}
    \item \textbf{Warstwa detonacji i weryfikacji kodu (\textit{Detonation Sandbox}):}
    \begin{itemize}
        \item \texttt{disp-eval-sandbox} -- jednorazowa maszyna wirtualna bez jakiegokolwiek interfejsu sieciowego (\texttt{NetVM = none}), powoływana wyłącznie na czas testów niezweryfikowanego kodu studenckiego.
    \end{itemize}
\end{enumerate}

\begin{figure}[htbp]
\centering
\resizebox{\linewidth}{!}{%
\begin{tikzpicture}[
    node distance=1.1cm and 1.2cm,
    every node/.style={font=\sffamily\footnotesize},
    box/.style={rectangle, draw=black!70, rounded corners=3pt, thick, align=center, fill=white, drop shadow},
    vault/.style={box, fill=purple!15, draw=purple!80},
    trusted/.style={box, fill=blue!15, draw=blue!80},
    osint/.style={box, fill=green!15, draw=green!80},
    untrusted/.style={box, fill=red!15, draw=red!80},
    infra/.style={box, fill=gray!15, draw=gray!80},
    xen/.style={rectangle, draw=black, thick, fill=yellow!20, minimum width=13cm, minimum height=0.9cm, font=\sffamily\bfseries\small},
    arrow/.style={-{Stealth[length=2.5mm]}, thick},
    dashedarrow/.style={-{Stealth[length=2.5mm]}, thick, dashed}
]

% Hypervisor Xen
\node[xen] (xen) at (0,0) {Hipernadzorca Xen 4.17 (Bare-Metal Type-1 Hypervisor)};

% Hardware
\node[box, fill=gray!30, below=0.5cm of xen, minimum width=13cm] (hw) {Fizyczny sprzęt: CPU (Intel VT-x / AMD-V), IOMMU (VT-d / AMD-Vi), TPM 2.0, Pamięć RAM};

% dom0
\node[box, fill=black!10, draw=black, above=0.7cm of xen.west, anchor=south west, minimum width=3.0cm, minimum height=1.6cm] (dom0) {\textbf{dom0} \\ (Zarządzanie GUI, \\ Qrexec Policy, \\ \textit{Brak sieci})};

% Domains
\node[vault, above=0.7cm of xen, xshift=-1.6cm, minimum width=2.4cm, minimum height=1.6cm] (vault) {\textbf{vault-sec} \\ Klucze SSH/GPG \\ Hasła \\ \textit{NetVM = none}};

\node[trusted, right=0.4cm of vault, minimum width=2.4cm, minimum height=1.6cm] (work) {\textbf{work-academic} \\ Dydaktyka, Publikacje \\ Dziekanat / Poczta};

\node[osint, right=0.4cm of work, minimum width=2.4cm, minimum height=1.6cm] (osint) {\textbf{disp-anon-osint} \\ Badania OSINT \\ Rynek pracy, granty \\ \textit{Whonix / Tor}};

\node[untrusted, right=0.4cm of osint, minimum width=2.5cm, minimum height=1.6cm] (deton) {\textbf{disp-eval-sandbox} \\ Detonacja kodu \\ Prace studentów \\ \textit{NetVM = none}};

% Infra domains below
\node[infra, above=2.7cm of work, xshift=-0.5cm, minimum width=2.2cm] (firewall) {\texttt{sys-firewall}};
\node[infra, left=0.4cm of firewall, minimum width=2.0cm] (net) {\texttt{sys-net} \\ (Karty NIC)};
\node[infra, right=0.4cm of firewall, minimum width=2.2cm] (whonix) {\texttt{sys-whonix} \\ (Brama Tor)};

% Connections
\draw[arrow] (work) -- (firewall);
\draw[arrow] (firewall) -- (net);
\draw[arrow] (osint) -- (whonix);
\draw[arrow] (whonix) -- (firewall);

% Qrexec channels
\draw[dashedarrow, purple!80] (work) -- node[above, sloped] {\tiny Split-SSH} (vault);
\draw[dashedarrow, red!80] (work) -- node[above, sloped] {\tiny qrexec ZT-SCEP} (deton);

\end{tikzpicture}%
}
\caption{Architektura modelu ACF (Academic Compartmentalization Framework) na bazie Qubes OS 4.2. Linie ciągłe oznaczają ruch sieciowy; linie przerywane wskazują bezpieczne wywołania procedur Qubes RPC (\texttt{qrexec}).}
\label{fig:acf_architecture}
\end{figure}

\subsection{Magistrala Qubes RPC (\texttt{qrexec}) jako szkielet komunikacji Zero-Trust}
Kluczowym elementem wyróżniającym Qubes OS na tle innych rozwiązań jest mechanizm komunikacji między-domenowej \textbf{Qubes RPC} (\texttt{qrexec})~\cite{qubes_qrexec2024}. W tradycyjnych systemach wymiana danych pomiędzy maszynami wirtualnymi wymaga utworzenia wirtualnego interfejsu sieciowego (vNIC), co natychmiast otwiera przestrzeń na ataki sieciowe (skanowanie portów, exploity stosu TCP/IP).

W Qubes OS \texttt{qrexec} opiera się na wydzielonych kanałach pamięci współdzielonej (\textit{vchan}) zarządzanych bezpośrednio przez hipernadzorcę Xen. Każde żądanie wywołania usługi RPC podlega rygorystycznej ewaluacji w silniku polityk znajdującym się w \texttt{dom0}. W systemie Qubes OS 4.2 wprowadzono modularną strukturę reguł zapisaną w katalogu \texttt{/etc/qubes/policy.d/}. Przykładowa reguła zezwalająca wyłącznie domenie programistycznej na odpytanie magazynu kluczy o podpis kryptograficzny bez udostępniania samego klucza prywatnego ma postać:
\begin{lstlisting}[language=bash, caption={Konfiguracja polityki Qubes RPC dla usługi Split-GPG w systemie Qubes OS 4.2.}]
# /etc/qubes/policy.d/30-academic-crypto.policy
qubes.Gpg  *  dev-trusted  vault-sec  ask  default_target=vault-sec
qubes.Gpg  *  *            vault-sec  deny
\end{lstlisting}

Powyższa polityka gwarantuje, że żadna inna domena (a w szczególności zanieczyszczone środowisko detonacji czy przeglądania stron) nie ma możliwości zażądania operacji kryptograficznej, a każda próba użycia klucza w \texttt{dev-trusted} wymaga jawnej akceptacji użytkownika w bezpiecznym oknie dialogowym generowanym bezpośrednio przez \texttt{dom0}.

\section{Scenariusze wdrożeniowe w działalności akademickiej}
\label{sec:scenariusze_akademickie}

Aby zweryfikować praktyczną użyteczność modelu ACF, zdefiniowano i zaimplementowano trzy kluczowe scenariusze operacyjne odpowiadające profilowi pracy studentów i wykładowców kierunków technicznych.

\subsection{Scenariusz 1: Bezpieczny wywiad jawnoźródłowy (OSINT)}
W procesie badawczym studenci i naukowcy poszukują literatury, analizują trendy rynkowe na portalach rekrutacyjnych, badają podatności publicznych systemów IT oraz weryfikują materiały nadsyłane ze źródeł otwartych. Działania te wiążą się z dwoma zasadniczymi zagrożeniami:
\begin{itemize}
    \item \textbf{Profilowanie behawioralne i fingerprinting:} przeglądanie komercyjnych serwisów badawczych i forów z poziomu standardowej przeglądarki umożliwia platformom analitycznym powiązanie zapytań z adresem IP uczelni oraz unikalnym odciskiem cyfrowym stacji roboczej;
    \item \textbf{Infekcja złośliwym dokumentem:} pobieranie plików PDF lub DOCX z nieoficjalnych repozytoriów grozi wykonaniem wrogiego kodu przy użyciu podatności w parserach dokumentów.
\end{itemize}

W ramach modelu ACF zaimplementowano dwuetapowy mechanizm ochrony:
\begin{enumerate}
    \item \textbf{Efemeryczna sesja Whonix (\texttt{disp-anon-whonix}):} przeglądanie zasobów realizowane jest w jednorazowej maszynie spiętej z bramą Tor. Wszystkie parametry stosu TCP, zegara systemowego oraz rozdzielczości ekranu podlegają unifikacji. Po zakończeniu wyszukiwania maszyna jest bezpowrotnie niszczona wraz ze wszystkimi ciasteczkami sesyjnymi i historią przeglądania.
    \item \textbf{Oczyszczanie dokumentów (\textit{Convert to Trusted PDF}):} pobrany z sieci niepewny plik PDF jest przesyłany z domeny OSINT do zaufanej domeny \texttt{work-academic} wyłącznie po przejściu procedury oczyszczania:
\begin{lstlisting}[language=bash, caption={Uruchomienie konwersji niepewnego dokumentu w bezpieczną bitmapę.}]
# Wywolanie z poziomu domeny OSINT do jednorazowego kontenera graficznego
qvm-convert-pdf publikacja_z_nieznanego_zrodla.pdf
\end{lstlisting}
Procedura ta w locie powołuje jednorazową maszynę wirtualną pozbawioną sieci, renderuje każdą stronę dokumentu do postaci surowej tablicy pikseli (RGB bitmap), a następnie składa na ich podstawie nowy dokument PDF. Wszelkie exploity binarne, złośliwe makra czy skrypty JavaScript osadzone w oryginalnym pliku zostają bezpowrotnie zneutralizowane.
\end{enumerate}

\subsection{Scenariusz 2: Wytwarzanie oprogramowania a ataki na łańcuch dostaw (Supply Chain)}
Współczesna dydaktyka programowania wymaga częstego korzystania z zewnętrznych menedżerów zależności (\texttt{npm}, \texttt{pip}, \texttt{cargo}). Jak wykazano w badaniach~\cite{pastrana2023malware}, napastnicy regularnie publikują w tych repozytoriach pakiety podszywające się pod popularne biblioteki (\textit{typosquatting}). Standardowy skrypt instalacyjny (np. \texttt{postinstall} w ekosystemie Node.js) uruchamia się z uprawnieniami bieżącego użytkownika i w tradycyjnym systemie może natychmiast odczytać katalogi \texttt{\textasciitilde/.ssh} czy \texttt{\textasciitilde/.aws/credentials}.

W systemie Qubes OS zastosowano mechanizm \textbf{Split-SSH}:
\begin{itemize}
    \item Klucze kryptograficzne są generowane i przechowywane w odciętej od sieci domenie \texttt{vault-sec};
    \item W domenie deweloperskiej \texttt{dev-trusted} klient \texttt{ssh} komunikuje się z atrapą gniazda (\textit{dummy socket}), która przekazuje żądania podpisania pakietu przez magistralę \texttt{qrexec};
    \item Daemon w \texttt{dom0} wyświetla natywne okno autoryzacji z pytaniem: \textit{„Czy zezwolić dev-trusted na dostęp do klucza id\_ed25519 w vault-sec?”}.
\end{itemize}
Nawet w przypadku instalacji skrajnie złośliwego pakietu z biblioteki PyPI/npm, atakujący nie ma fizycznej możliwości skopiowania klucza prywatnego, ponieważ klucz ten nigdy nie pojawia się w pamięci RAM maszyny programistycznej.

\subsection{Scenariusz 3: Zautomatyzowany potok bezpiecznej detonacji i weryfikacji prac studenckich (ZT-SCEP)}
Weryfikacja prac laboratoryjnych z przedmiotów takich jak programowanie niskopoziomowe, algorytmy czy technologie sieciowe wiąże się z koniecznością kompilacji i uruchomienia setek plików nadesłanych przez studentów. Kod ten często zawiera:
\begin{itemize}
    \item błędy skutkujące zawieszeniem stacji (przypadkowe pętle \textit{fork bomb}, nieskończona alokacja pamięci);
    \item celowe próby naruszenia środowiska oceniania (np. skanowanie portów w sieci lokalnej, próby odczytania testów sprawdzających, otwieranie gniazd nasłuchujących \textit{reverse shell}).
\end{itemize}

W ramach niniejszych badań opracowano potok \textbf{Zero-Trust Student Code Evaluation Pipeline (ZT-SCEP)}. Przepływ ewaluacji realizowany jest w następujących krokach:
\begin{enumerate}
    \item Prace studenckie (zarchiwizowane w formacie \texttt{.zip} lub \texttt{.tar.gz}) są gromadzone w domenie dydaktycznej \texttt{work-academic};
    \item Skrypt weryfikujący wywołuje nowo powołaną maszynę jednorazową (\texttt{DispVM}) bazującą na szablonie \texttt{tpl-eval-sandbox}, którego interfejs sieciowy jest trwale wyłączony (\texttt{netvm=""});
    \item Archiwum z kodem jest strumieniowane przez \texttt{qrexec} bezpośrednio na standardowe wejście (\texttt{stdin}) procesu ewaluacyjnego w maszynie efemerycznej;
    \item Skrypt ewaluatora w \texttt{DispVM} nakłada rygorystyczne limity zasobów poprzez podsystem \texttt{systemd-run} / \texttt{cgroups} (maksymalnie 512~MB RAM, 10 sekund czasu CPU, zakaz tworzenia procesów potomnych powyżej 32 wątków);
    \item Wyniki testów jednostkowych, raport z kompilacji oraz kod wyjścia są odsyłane przez \texttt{stdout} do domeny dydaktycznej;
    \item Hipernadzorca Xen niszczy maszynę wirtualną wraz z całym stanem pamięci i dysku ulotnego.
\end{enumerate}

\begin{lstlisting}[language=bash, caption={Skrypt zarządcy potoku ZT-SCEP w domenie dydaktycznej (\texttt{work-academic}).}]
#!/usr/bin/env bash
# run_ztscep.sh -- uruchomienie bezpiecznej ewaluacji kodu studenckiego
SUBMISSION_ARCHIVE="$1"
STUDENT_ID="$2"

echo "[*] Inicjalizacja izolowanej domeny detonacyjnej (DispVM)..."
# Wywolanie uslugi RPC z przekierowaniem strumienia archiwum na stdin
qrexec-client-vm @dispvm:tpl-eval-sandbox \
    academic.EvaluateSubmission < "${SUBMISSION_ARCHIVE}" \
    > "raport_${STUDENT_ID}.json"

EXIT_STATUS=$?
if [ $EXIT_STATUS -eq 0 ]; then
    echo "[+] Ocena zakonczona sukcesem. Raport zapisany w work-academic."
else
    echo "[-] Wykryto anomalie lub przekroczenie limitow w kodzie studenta!"
fi
\end{lstlisting}

Dzięki uniezależnieniu potoku od stosu sieciowego oraz wykorzystaniu jednorazowych maszyn Xen, stacja robocza wykładowcy pozostaje w 100\% odizolowana od jakichkolwiek skutków ubocznych badanego kodu.

\section{Ewaluacja empiryczna i analiza bezpieczeństwa}
\label{sec:ewaluacja_empiryczna}

W celu ilościowej i jakościowej oceny proponowanego modelu ACF przeprowadzono serię testów wydajnościowych oraz eksperymentów penetracyjnych w dedykowanym środowisku laboratoryjnym.

\subsection{Stanowisko badawcze i metodologia pomiarów}
Eksperymenty zrealizowano na stacji roboczej o konfiguracji typowej dla współczesnych laboratoriów akademickich:
\begin{itemize}
    \item Procesor: AMD Ryzen 7 7840HS (8 rdzeni, 16 wątków, obsługa AMD-V oraz IOMMU/AMD-Vi);
    \item Pamięć operacyjna: 32~GB DDR5-5600;
    \item Pamięć masowa: szybki dysk SSD NVMe PCIe 4.0 o przepustowości 5000~MB/s;
    \item Oprogramowanie bazowe: Qubes OS 4.2 z jądrem Linux 6.6 i hipernadzorcą Xen 4.17.
\end{itemize}

Pomiary czasu powoływania maszyn efemerycznych oraz przepustowości wykonano w 100 powtórzeniach, rejestrując wartości średnie oraz odchylenie standardowe.

\subsection{Wyniki wydajnościowe potoku ZT-SCEP}
W pracy dydaktycznej kluczowym parametrem ergonomii jest narzut czasowy związany z powoływaniem środowiska weryfikacji. Prowadzący sprawdzający 30 prac laboratoryjnych nie może oczekiwać kilkudziesięciu sekund na start każdej maszyny wirtualnej.

\begin{table}[htbp]
\centering
\caption{Porównanie parametrów rozruchu i narzutu pamięciowego środowisk testowych.}
\label{tab:performance_comparison}
\resizebox{\linewidth}{!}{%
\begin{tabular}{@{}lcccc@{}}
\toprule
\textbf{Rozwiązanie} & \textbf{Typ izolacji} & \textbf{Czas startu [s]} & \textbf{Zużycie RAM [MB]} & \textbf{Twarda bariera} \\
\midrule
Tradycyjny host (Ubuntu) & Brak (Proces lokalny) & $0{,}05 \pm 0{,}01$ & 15 & NIE \\
Kontener Docker & Współdzielone jądro & $0{,}42 \pm 0{,}05$ & 45 & NIE (Wspólny kernel) \\
VirtualBox VM & Hipernadzorca typu 2 & $14{,}80 \pm 1{,}20$ & 2048 & CZĘŚCIOWO \\
\textbf{Qubes DispVM (ACF)} & \textbf{Hipernadzorca Xen} & $\mathbf{2{,}42 \pm 0{,}18}$ & \textbf{450} & \textbf{TAK (Xen Type-1)} \\
\bottomrule
\end{tabular}%
}
\end{table}

Jak wynika z tabeli~\ref{tab:performance_comparison}, czas powołania jednorazowej maszyny \texttt{DispVM} w architekturze Qubes OS wynosi średnio \textbf{2,42 sekundy}, przy narzucie pamięci operacyjnej na poziomie zaledwie 450~MB (dzięki dynamicznemu mechanizmowi Xen Memory Ballooning oraz współdzieleniu szablonu bazowego w trybie tylko do odczytu). Całkowite zniszczenie maszyny i zwolnienie zasobów pamięciowych zajmuje poniżej \textbf{0,35 sekundy}.

Przepustowość przesyłania archiwów studenckich przez magistralę pamięci współdzielonej \texttt{vchan} w usłudze Qubes RPC wyniosła średnio \textbf{482 MB/s}. Oznacza to, że przesłanie standardowego archiwum z projektem o rozmiarze 5--15~MB trwało mniej niż 35 milisekund, stanowiąc pomijalny ułamek całego cyklu testowego.

\subsection{Odporność na celowe i przypadkowe zagrożenia w kodzie studenckim}
W celu zbadania szczelności środowiska ZT-SCEP przeprowadzono cztery syntetyczne testy penetracyjne, symulujące najczęstsze wektory awarii oraz celowych naruszeń bezpieczeństwa:

\begin{enumerate}
    \item \textbf{Wyczerpanie zasobów (Pętla widelcowa -- \textit{Fork Bomb}):} uruchomiono rekurencyjny proces generujący procesy potomne. Nałożone wewnątrz szablonu \texttt{DispVM} limity cgroups (\texttt{pids.max = 32}) skutecznie zablokowały tworzenie kolejnych wątków. Co kluczowe, system gospodarza (\texttt{dom0}) oraz domena dydaktyczna \texttt{work-academic} nie odnotowały jakiegokolwiek spadku responsywności (obciążenie CPU w \texttt{dom0} pozostało poniżej 1,8\%).
    \item \textbf{Próba eksfiltracji danych przez gniazdo sieciowe (\textit{Reverse Shell}):} kod studenta próbował wysłać zawartość katalogu bieżącego na zewnętrzny adres IP za pomocą gniazda TCP. Ze względu na sztywną konfigurację \texttt{NetVM = none}, próba wywołania funkcji systemowej \texttt{connect()} zakończyła się natychmiastowym błędem \texttt{ENETUNREACH} (Network is unreachable). Wyciek danych poza maszynę wirtualną był fizycznie niemożliwy.
    \item \textbf{Złośliwe skrypty instalacyjne (\textit{Supply Chain Typosquatting}):} w projekcie zasymulowano złośliwy hook \texttt{preinstall} próbujący skopiować klucze \texttt{\textasciitilde/.ssh/id\_ed25519}. Ponieważ katalog domowy w czystej instancji \texttt{DispVM} nie zawierał żadnych tożsamości, a usługa Split-SSH nie była powiązana z domeną testową, atak zakończył się niepowodzeniem.
    \item \textbf{Destrukcja systemu plików (\texttt{rm -rf /}):} kod wykonywał próbę destrukcji partycji systemowej z uprawnieniami superużytkownika (\texttt{root}). Zmiany objęły wyłącznie ulotną warstwę różnicową pamięci podręcznej (\textit{ephemeral snapshot}). Po wyjściu z procedury RPC i anihilacji maszyny wirtualnej, szablon bazowy pozostał całkowicie nienaruszony.
\end{enumerate}

\begin{table}[htbp]
\centering
\caption{Porównanie odporności środowiska ewaluacji na wybrane scenariusze zagrożeń.}
\label{tab:security_comparison}
\resizebox{\linewidth}{!}{%
\begin{tabular}{@{}lccc@{}}
\toprule
\textbf{Wektor ataku} & \textbf{Docker} & \textbf{VirtualBox} & \textbf{Qubes ACF (ZT-SCEP)} \\
\midrule
Fork Bomb na hoście & Ryzyko (DoW hosta) & Izolacja procesu & \textbf{Pełna izolacja sprzętowa} \\
Eksfiltracja sieciowa & Wymaga reguł bridge & Wymaga odpięcia kabla & \textbf{Fizyczny brak vNIC} \\
Kradzież tożsamości SSH & Możliwa przy bind mount & Możliwa przez schowek & \textbf{Klucze w odciętej domenie} \\
Zniszczenie systemu & Ryzyko kontenera & Manualny revert & \textbf{Automatyczna anihilacja VM} \\
\bottomrule
\end{tabular}%
}
\end{table}

Uzyskane wyniki jednoznacznie dowodzą, że model ACF zapewnia bezprecedensowy poziom odporności na błędy i ataki, przewyższając powszechnie stosowane kontenery oraz maszyny wirtualne typu 2.

\section{Dyskusja, ergonomia i aspekty wdrożeniowe}
\label{sec:dyskusja_i_wdrozenie}

Mimo bezdyskusyjnych zalet w zakresie bezpieczeństwa, wdrożenie architektury opartej na kompartmentaryzacji w środowisku akademickim napotyka na szereg wyzwań o charakterze technologicznym, ergonomicznym i organizacyjnym.

\subsection{Krzywa uczenia i bariera wejścia}
Najczęściej podnoszonym argumentem przeciwko powszechnemu stosowaniu Qubes OS jest stroma krzywa uczenia (\textit{steep learning curve}) oraz konieczność redefinicji przyzwyczajeń użytkownika:
\begin{itemize}
    \item \textbf{Higiena mentalna i identyfikacja domen:} użytkownik musi wypracować nawyk operowania w strefach zaufania. Kluczowym ułatwieniem w Qubes OS jest kolorystyczne kodowanie obramowań okien (np. czerwone dla niezweryfikowanych stron WWW, niebieskie dla spraw uczelnianych, czarne dla odciętego magazynu haseł). Badania wykazują, że po okresie adaptacji (trwającym zwykle 1--2 tygodnie) wskaźnik pomyłkowego przeklejania wrażliwych danych spada niemal do zera;
    \item \textbf{Izolacja schowka (\textit{Inter-VM Clipboard}):} Qubes uniemożliwia bezpośrednie współdzielenie schowka pomiędzy domenami. Kopiowanie danych wymaga dwuetapowej sekwencji klawiszy (\texttt{Ctrl+C} w oknie źródłowym, a następnie \texttt{Ctrl+Shift+C} do globalnego bufora \texttt{dom0}), co chroni przed atakami typu \textit{clipboard hijacking}, lecz wymaga zdyscyplinowania użytkownika;
    \item \textbf{Wymagania sprzętowe:} efektywne działanie hipernadzorcy Xen wymaga stacji roboczej wyposażonej w co najmniej 16~GB pamięci RAM (rekomendowane 32~GB) oraz szybki dysk NVMe. Ze względu na bezpieczeństwo GUI, Qubes OS celowo nie wspiera natywnej akceleracji GPU w domenach nieuprzywilejowanych. Choć stanowi to ograniczenie w modelowaniu 3D czy lokalnym trenowaniu sieci neuronowych, zadania te w warunkach akademickich są zazwyczaj delegowane na dedykowane serwery obliczeniowe zarządzane centralnie (np. klastry Proxmox~VE~\cite{zarnecki2026llm_proxmox}).
\end{itemize}

\subsection{Zgodność z RODO i suwerenność cyfrowa instytucji}
W dobie zaostrzonych regulacji dotyczących ochrony danych osobowych (Rozporządzenie RODO / GDPR), uczelnie wyższe ponoszą bezpośrednią odpowiedzialność prawną za poufność przetwarzanych informacji. Wykładowcy przechowują na stacjach roboczych wrażliwe dane: numery albumów, wyniki egzaminów, opinie promotorskie oraz dane teleadresowe studentów.

Model ACF bezpośrednio realizuje kluczowe zasady prawa ochrony danych:
\begin{enumerate}
    \item \textbf{Zasada integralności i poufności (art. 5 ust. 1 lit. f RODO):} separacja bazy ocen w domenie \texttt{work-academic} uniemożliwia jej kradzież przez złośliwy kod studencki detonowany w domenie \texttt{DispVM};
    \item \textbf{Zasada minimalizacji i rozliczalności (art. 5 ust. 1 lit. c RODO):} wykorzystanie bezśladowego potoku OSINT opartego na maszynach jednorazowych zapobiega gromadzeniu niepotrzebnych danych telemetrycznych i ciasteczek śledzących;
    \item \textbf{Suwerenność technologiczna:} Qubes OS jest projektem w pełni otwartoźródłowym, poddawanym ciągłemu audytowi społeczności międzynarodowej. W przeciwieństwie do komercyjnych systemów operacyjnych z telemetrią chmurową, żadne metadane dotyczące aktywności naukowej czy dydaktycznej nie są przesyłane do zewnętrznych korporacji technologicznych.
\end{enumerate}

\subsection{Rekomendacje metodyczne dla jednostek akademickich}
Aby obniżyć barierę adaptacji i umożliwić studentom oraz pracownikom płynne korzystanie z zalet kompartmentaryzacji, zaleca się:
\begin{itemize}
    \item Przygotowanie przez uczelniany dział IT ustandaryzowanych szablonów (\textit{golden templates}) maszyn wirtualnych z preinstalowanymi narzędziami programistycznymi i dydaktycznymi;
    \item Udostępnienie predefiniowanego zestawu reguł Qubes RPC w pliku konfiguracyjnym polityki (np. \texttt{30-academic.policy}), co eliminuje konieczność ręcznej konfiguracji mechanizmów Split-SSH i ZT-SCEP przez każdego użytkownika;
    \item Wdrożenie warsztatów z zakresu higieny operacyjnej (OpSec) na wczesnych semestrach studiów informatycznych, traktujących pracę w środowisku kompartmentaryzowanym jako element kanonu nowoczesnej inżynierii bezpieczeństwa.
\end{itemize}

\section{Podsumowanie i wnioski}
\label{sec:podsumowanie}

W dobie dynamicznie ewoluujących cyberzagrożeń, paradygmat ochrony stacji roboczych oparty na tradycyjnych systemach monolitycznych i antywirusach stał się dalece niewystarczający w specyficznym, otwartym środowisku akademickim. Konieczność codziennego przetwarzania niezweryfikowanych materiałów z sieci, poszukiwania informacji jawnoźródłowych (OSINT), instalowania setek zewnętrznych zależności programistycznych oraz laboratoryjnego uruchamiania kodu nadesłanego przez studentów wymaga wdrożenia bezwzględnej izolacji sprzętowej.

W niniejszym artykule przedstawiono kompleksową koncepcję, architekturę oraz ewaluację empiryczną modelu \textbf{Academic Compartmentalization Framework (ACF)}, opartego na systemie \textbf{Qubes OS 4.2} i hipernadzorcy typu 1 \textbf{Xen}. Do najważniejszych osiągnięć pracy należą:
\begin{enumerate}
    \item \textbf{Opracowanie bezśladowego potoku badawczego OSINT:} integracja efemerycznych instancji Whonix z mechanizmem automatycznego oczyszczania dokumentów (\textit{Convert to Trusted PDF}) umożliwiła prowadzenie zaawansowanego wywiadu otwartego bez ryzyka profilowania cyfrowego, ujawnienia tożsamości sieciowej uczelni czy infekcji złośliwymi exploitami;
    \item \textbf{Wdrożenie ochrony tożsamości kryptograficznej programisty:} zastosowanie architektury Split-SSH i Split-GPG pozwoliło na całkowite unieszkodliwienie wektorów kradzieży poświadczeń wynikających z ataków typu \textit{supply chain} w ekosystemach pakietów \texttt{npm} i \texttt{pip};
    \item \textbf{Stworzenie i przetestowanie potoku ZT-SCEP:} autorski mechanizm ewaluacji prac studenckich z wykorzystaniem magistrali Qubes RPC (\texttt{qrexec}) i jednorazowych maszyn wirtualnych bez interfejsu sieciowego (\texttt{DispVM}) zagwarantował 100\% odporność stacji roboczej na celowe i przypadkowe exploity (w tym pętle \textit{fork bomb}, próby nawiązania \textit{reverse shell} oraz destrukcję plików);
    \item \textbf{Weryfikacja wydajnościowa i użytkowa:} pomiary empiryczne potwierdziły, że czas powołania maszyny jednorazowej (średnio 2,42~s) oraz narzut pamięciowy (450~MB) w pełni predysponują to rozwiązanie do zastosowań w codziennej praktyce laboratoryjnej.
\end{enumerate}

Kierunkiem dalszych prac badawczo-rozwojowych będzie opracowanie zautomatyzowanego mostu integrującego potok ZT-SCEP z uczelnianymi systemami zarządzania nauczaniem (LMS, takimi jak Moodle czy Canvas) za pośrednictwem asynchronicznych brokerów komunikatów, a także integracja fizycznych tokenów sprzętowych FIDO2 z domeną kluczy \texttt{vault-sec} w celu zapewnienia pełnego łańcucha zaufania od warstwy fizycznej po dydaktyczną.


%%----------------------------------------------
\begin{singlespace}
\printbibliography[title={Bibliografia}]
\end{singlespace}

\end{document}
